\item A lo largo del tiempo han existido diversos modelos de datos, como el modelo jerárquico, el modelo de red, el modelo relacional, el modelo orientado a objetos y los modelos NoSQL (llave-valor, documento, columnares, grafos). Responde lo siguiente: Elabora una tabla comparativa entre al menos cuatro modelos de datos, destacando sus principales características, ventajas, desventajas y casos de uso típicos. ¿Por qué crees que el modelo relacional se convirtió en el más ampliamente adoptado durante décadas? ¿En qué contextos actuales los modelos NoSQL superan al modelo relacional en eficiencia o flexibilidad? Justifica con ejemplos.